With the rapid development of the Internet, email has become a fast and convenient method of communication that is widely used in our daily lives. 

However, the popularity of email has also caused a problem that cannot be ignored: spam. Unscrupulous businessmen and criminals have exploited this communication medium to send spam with commercial advertisements and malicious or illegal files to users on a large scale \cite {kumar_email_2020}. Those spam emails not only occupy network resources and cause network congestion, but also affect people's normal work and life.

Therefore, how to quickly and accurately check whether an email is spam or not 
becomes a major topic worthy of research. It can protect people's property and 
privacy security and also improve people's work and life efficiency. 

In this paper, we will try to solve the spam email detection problem using the latest transformer based models, and using the TREC 2007 spam email dataset to train our model \cite{trecspam2007}.
Considering the growing use of pre-trained language in the NLP systems, in this experiment, we will adopt two distinct transformer based models, BERT(encoder-only model) \cite{devlin2018bert} and GPT-2(decoder-only model) \cite{radford2019language}, for this spam detection task, and to verify their real performance on our task. 
This comparative study will also try to find out if the generative pretraining (GPT-2) model can provide any advantages or disadvantages in a spam detection problem.